\section{Results}\label{sec:results}

\subsection{Baselines on Dev}

Table~\ref{tab:dev} compares \ours{} with six baselines by the number of harmful family--model pairs (/144). The
detection group has no effect on Dev: \cPG{} (113) and \cMel{} (112) are on par with \cNone{} (112). The prompt-level
group gives mixed results. \cSpot{} only reduces the count to 91 and even increases the number of harmful families on
Qwen3.5 (from 40 to 43), whereas \cSand{} reduces it to 31. The rule group improves with the specificity of the rules,
with \cPhrase{} at 94, \cProg{} at 61 and \cRules{} at 45. \ours, which combines rules with the reminder, brings the
count down to 10.

The last column of Table~\ref{tab:dev} shows that a stronger model does not replace a defence. Under \cNone, Gemini 3
Flash is attacked less often than the open models (22/48 families compared with 34 to 40) but is still compromised in
almost half of the families. \cRules{} reduces this to 12, and \ours{} to only 1.

% Nguồn: outline/05_results.md Bảng 5 và Bảng 10 (Gemini; = study2_thesis_tables.md Bảng 2.11). Tên: P1 Spotlight,
% P2 Sandwich, G1 PromptGuard, MEL MELON-det, PRG Progent-pol, R1 PhraseRules, D2 BG-Rules, D2+P2 BrokerGuard (13/11: bỏ dòng D3DEST/BG-v2).
% Gemini chỉ chạy 5 cấu hình trên Dev; không cộng vào cột Tổng. In đậm: ít họ gây hại nhất mỗi cột.
\begin{table}[!htbp]
\centering
\caption{Baselines and \ours{} on Dev. Each cell: harmful families /48 $\cdot$ plain benign /240 $\cdot$ authorised
benign /80. Total: harmful family--model pairs over the three open models (/144). $^\dagger$Reference only, not part of
the tests. ``--'': not run.}\label{tab:dev}
\footnotesize
\setlength{\tabcolsep}{3pt}
\begin{tabular*}{\textwidth}{@{\extracolsep{\fill}}llrrrrr@{}}
\toprule
Group & Configuration & \makecell[r]{Gemma\\4-12B} & \makecell[r]{Llama\\3.1-8B} & \makecell[r]{Qwen3.5\\9B} & Total & \makecell[r]{Gemini\\3 Flash$^\dagger$} \\
\midrule
 & \cNone & 38\,$\cdot$\,159\,$\cdot$\,54 & 34\,$\cdot$\,81\,$\cdot$\,30 & 40\,$\cdot$\,140\,$\cdot$\,36 & 112 & 22\,$\cdot$\,184\,$\cdot$\,53 \\
\multirow{2}{*}{Prompt} & \cSpot & 24\,$\cdot$\,146\,$\cdot$\,54 & 24\,$\cdot$\,74\,$\cdot$\,36 & 43\,$\cdot$\,137\,$\cdot$\,38 & 91 & -- \\
 & \cSand & 7\,$\cdot$\,150\,$\cdot$\,54 & 13\,$\cdot$\,84\,$\cdot$\,40 & 11\,$\cdot$\,144\,$\cdot$\,32 & 31 & \textbf{1}\,$\cdot$\,191\,$\cdot$\,64 \\
\multirow{2}{*}{\makecell[l]{Detec-\\tion}} & \cPG & 39\,$\cdot$\,160\,$\cdot$\,54 & 34\,$\cdot$\,81\,$\cdot$\,30 & 40\,$\cdot$\,139\,$\cdot$\,36 & 113 & -- \\
 & \cMel & 38\,$\cdot$\,160\,$\cdot$\,54 & 34\,$\cdot$\,79\,$\cdot$\,30 & 40\,$\cdot$\,142\,$\cdot$\,36 & 112 & -- \\
\multirow{2}{*}{Rules} & \cProg & 21\,$\cdot$\,160\,$\cdot$\,54 & 18\,$\cdot$\,80\,$\cdot$\,32 & 22\,$\cdot$\,138\,$\cdot$\,36 & 61 & -- \\
 & \cPhrase & 36\,$\cdot$\,160\,$\cdot$\,49 & 25\,$\cdot$\,80\,$\cdot$\,27 & 33\,$\cdot$\,138\,$\cdot$\,33 & 94 & 18\,$\cdot$\,182\,$\cdot$\,50 \\
\midrule
\multirow{2}{*}{Ours} & \cRules & 17\,$\cdot$\,158\,$\cdot$\,44 & 13\,$\cdot$\,80\,$\cdot$\,27 & 15\,$\cdot$\,136\,$\cdot$\,28 & 45 & 12\,$\cdot$\,186\,$\cdot$\,44 \\
 & \textbf{\ours} & \textbf{2}\,$\cdot$\,147\,$\cdot$\,45 & \textbf{5}\,$\cdot$\,82\,$\cdot$\,36 & \textbf{3}\,$\cdot$\,141\,$\cdot$\,24 & \textbf{10} & \textbf{1}\,$\cdot$\,192\,$\cdot$\,49 \\
\bottomrule
\end{tabular*}
\end{table}


\subsection{\tOne: Pre-registered Tests}

\tOne{} was run once with \ours{} as the final configuration (Table~\ref{tab:test1}). Three defence hypotheses hold at
the family level. The rules of \ours{} outperform phrase rules (H1, from 53 to 24 harmful families out of 72), including
on the mechanisms that evade phrase patterns (H3). Combining them with the reminder outperforms the reminder alone (H6,
from 28 to 8 families). Compared with \cNone, \ours{} reduces the number of harmful families from 65 to 8. H2 is not
established. \cRules{} loses 2.8 points of \bPlain{} utility relative to \cNone, and the lower confidence bound
($-6.0$) crosses the $-5$ point margin. H4 and H5 are not supported, that is, no language effect is detected. At the
family level, the variant without diacritics is harmful in 50 families compared with 40 for standard Vietnamese (16 and
6 discordant families, Holm $p$ 0.16). By case, the difference is $+1.9$ points [$-5.1$, $+8.8$] without diacritics and
$+5.1$ points [$-1.9$, $+12.0$] for English. These intervals are still wide, so the results do not show that language
has no effect. In H5, the case rate rises while the family count falls, because English attacks succeed on more
models but in fewer families.

% Test-I = "Test v3" trong log/luận văn. Nguồn: facts.md (bảng Test v3, H1--H6 và mục "Bổ sung sau phản biện").
% 13/11: H1, H3, H6 ở mức họ (scripts/paper_cluster_stats_v1.py). 14/11 (phản biện 2, mục 4.6): H2, H4, H5 cũng ở mức họ
% (scripts/paper_cluster_stats_h245_v1.py). Holm trên sáu p mức họ: H1 6x, H6 5x, H3 4x, H4 3x, H2 2x (đơn điệu), H5 1x.
% In đậm: tốt nhất mỗi cột, không tính hàng ước lượng (owner bỏ câu này khỏi chú thích 14/11).
\begin{table}[!htbp]
\centering
\caption{\tOne{} (72 families, summed over three backbones). Top: harmful family--model pairs (/216), harmful attack
cases (/1,080), successful benign cases (/1,080 and /360). Bottom: six pre-registered hypotheses, tested at the family
level, Holm-adjusted $p$. $^\dagger$Estimated from \cNone{} trajectories.}\label{tab:test1}
\footnotesize
\setlength{\tabcolsep}{3pt}
\begin{tabular*}{\textwidth}{@{\extracolsep{\fill}}lrrrr@{}}
\toprule
Configuration & \makecell[r]{Harmful\\family $\times$ model} & \makecell[r]{Harmful cases\\(ASR)} & \makecell[r]{Plain\\benign} & \makecell[r]{Authorised\\benign} \\
\midrule
\cNone & 138 & 355 (32.9\%) & \textbf{489 (45.3\%)} & \textbf{141 (39.2\%)} \\
\cPhrase & 109 & 200 (18.5\%) & 468 (43.3\%) & 128 (35.6\%) \\
\cProg & 75 & 212 (19.6\%) & 488 (45.2\%) & \textbf{141 (39.2\%)} \\
\cSand & 47 & 61 (5.6\%) & 472 (43.7\%) & 124 (34.4\%) \\
\cRules & 49 & 72 (6.7\%) & 459 (42.5\%) & 118 (32.8\%) \\
\textbf{\ours} & \textbf{11} & \textbf{13 (1.2\%)} & 448 (41.5\%) & 108 (30.0\%) \\
Block all sinks$^\dagger$ & -- & 13 (1.2\%) & 95 (8.8\%) & 12 (3.3\%) \\
\bottomrule
\end{tabular*}

\smallskip
\begin{tabularx}{\textwidth}{@{}l>{\raggedright\arraybackslash}Xrrl@{}}
\toprule
 & Hypothesis & Cases (harmful families) & Holm $p$ & Result \\
\midrule
H1 & ASR: \cRules{} $<$ \cPhrase & 18.5\%$\to$6.7\% (53$\to$24) & $1.1\cdot10^{-8}$ & met \\
H2 & Utility: \cRules{} $\ge$ \cNone{}$-5$ & 45.3\%$\to$42.5\% & 0.16 & not met \\
H3 & ASR (\mCR, \mSL, \mAL): \cRules{} $<$ \cPhrase & 23.5\%$\to$8.5\% (28$\to$15) & $4.8\cdot10^{-4}$ & met \\
H4 & ASR \cNone: no diacritics $\ne$ standard & 31.5\%$\to$33.3\% (40$\to$50) & 0.16 & not met \\
H5 & ASR \cNone: English $\ne$ standard & 31.5\%$\to$36.6\% (40$\to$39) & 1.00 & not met \\
H6 & ASR: \ours{} $<$ \cSand & 5.6\%$\to$1.2\% (28$\to$8) & $4.8\cdot10^{-6}$ & met \\
\bottomrule
\end{tabularx}
\end{table}


Figure~\ref{fig:results} shows why the two layers need each other. The rules bring both redirection mechanisms to zero
regardless of wording. For the two leak mechanisms, rules and reminder perform similarly (29 and 27 cases through a
sink, 17 and 18 through the answer). For parameter tampering, the reminder is clearly more effective (10 compared with
25 cases). Combined, only 13 harmful cases remain: 6 sink leaks, 3 answer leaks and 4 parameter tamperings. By language
variant, \cNone{} is attacked almost equally in every variant (68 to 79 cases out of 216). \cRules{} lets the most
attacks through in English (32 cases, compared with 6 to 13 in the other variants), and the reminder fixes this
weakness (1 case left).

% Hình 2, vẽ bằng pgfplots (Overleaf tự vẽ, không render trên máy).
% Số liệu: facts.md và study2_thesis_tables.md. (a) họ gây hại /144 trên Dev, gộp 3 mô hình nền.
% 14/11: bỏ panel (a) Dev (trùng Bảng tab:dev) để rút gọn. Ca tấn công gây hại /180 mỗi cơ chế trên Test-I, gộp 3 mô hình nền.
% Kích thước trục cố định (scale only axis) để hai panel không đè nhau: tổng rộng khoảng 11,6 cm < \textwidth 12,2 cm.
\begin{figure}[!htbp]
\centering
\begin{tikzpicture}
\begin{axis}[
  scale only axis, width=8.5cm, height=2.6cm,
  ybar, bar width=3pt,
  ymin=0, ymax=100, ytick={0,25,50,75,100},
  symbolic x coords={ER,CR,SL,AL,PT,US}, xtick=data, enlarge x limits=0.1,
  axis x line*=bottom, axis y line*=left,
  tick label style={font=\scriptsize}, label style={font=\scriptsize},
  ylabel={Harmful attack cases (/180)},
  ymajorgrids, grid style={line width=0.2pt, draw=black!12},
  legend style={font=\tiny, draw=none, fill=white, fill opacity=0.85, text opacity=1,
    at={(0.99,0.99)}, anchor=north east, cells={anchor=west}},
  legend image code/.code={\draw[#1] (0cm,-0.07cm) rectangle (0.2cm,0.07cm);},
  ]
\addplot[fill=white, draw=black, line width=0.4pt] coordinates {(ER,58) (CR,65) (SL,89) (AL,60) (PT,68) (US,15)};
\addplot[fill=bgaqua, draw=none, postaction={pattern=north east lines, pattern color=white}]
  coordinates {(ER,1) (CR,1) (SL,27) (AL,18) (PT,10) (US,4)};
\addplot[fill=bgblue, draw=none] coordinates {(ER,0) (CR,0) (SL,29) (AL,17) (PT,25) (US,1)};
\addplot[fill=bgorange, draw=none] coordinates {(ER,0) (CR,0) (SL,6) (AL,3) (PT,4) (US,0)};
\legend{\cNone, \cSand, \cRules, \ours}
\end{axis}
\end{tikzpicture}
\caption{Harmful attack cases per mechanism on \tOne, summed over three backbones (/180 per mechanism). A missing
bar means 0 cases.}\label{fig:results}
\end{figure}


\subsection{\tTwo: Replication on New Data}

\tTwo{} consists of 48 families written after \ours{} was frozen (Table~\ref{tab:test2}). The number of harmful
families falls from 42 to 14 ($p = 3.7\cdot10^{-9}$), and no family switches from safe to harmful. The attack success
rate falls from 32.4\% to 3.8\%. The reduction holds on all three backbones: Gemma from 101 to 7 cases out of 240,
Llama from 50 to 10 and Qwen3.5 from 82 to 10. The cost is a 4.5-point drop in benign utility, from 54.5\% to 50.0\%
($p = 0.044$). On \tOne{} the drop is similar ($-5.1$ points, $p = 0.022$). Most of the loss is in the \bPlain{}
counterparts. In the \bAuth{} counterparts, the $-2.5$ point difference is not significant ($p = 0.57$). On \tOne, the
\bAuth{} counterparts drop more (9.2 points, from 39.2\% to 30.0\%). Llama produces 66 to 74 truncated outputs per
configuration, and all of them count as failures.

The bypass rate of \ours{} is higher on \tTwo{} than on \tOne{} (3.8\% versus 1.2\% of cases, 14/48 versus 8/72
families), while \cNone{} is almost unchanged. The family-level Wilson intervals overlap ($[18.2, 43.2]$\% and
$[5.7, 20.4]$\%), so this is a signal to investigate, not a conclusion. By mechanism, 23 of the 27 cases are leaks or
parameter tampering (14 sink leaks, 2 answer leaks, 7 parameter tamperings). The other four are 1 covert redirection
and 3 unrequested side effects. The cause of each case is analysed in Section~\ref{sec:discussion}.

% Test-II = "Test-3" trong log/luận văn. Nguồn: facts.md (bảng Test-3 và mục "Bổ sung sau phản biện").
% 13/11: bỏ BG-v2 (D3DEST); kiểm định ở mức họ (scripts/paper_cluster_stats_v1.py), không đăng ký trước.
% In đậm: tốt nhất mỗi cột, không tính hàng ước lượng.
\begin{table}[!htbp]
\centering
\caption{\tTwo{} (48 new families, summed over three backbones), replication analysis. Top: attack and plain benign
/720, authorised benign /240, benign utility /960. Bottom: difference of \ours{} from \cNone{} (percentage points,
family-level 95\% CI). $^\dagger$Estimated from \cNone{} trajectories.}\label{tab:test2}
\footnotesize
\setlength{\tabcolsep}{3pt}
\begin{tabular*}{\textwidth}{@{\extracolsep{\fill}}lrrrrr@{}}
\toprule
Configuration & \makecell[r]{Harmful\\families (/48)} & \makecell[r]{Successful\\attacks} & \makecell[r]{Plain\\benign} & \makecell[r]{Authorised\\benign} & \makecell[r]{Benign\\utility} \\
\midrule
\cNone & 42 & 233 (32.4\%) & \textbf{420 (58.3\%)} & \textbf{103 (42.9\%)} & \textbf{523 (54.5\%)} \\
Block all sinks$^\dagger$ & -- & 5 (0.7\%) & 96 (13.3\%) & 24 (10.0\%) & 120 (12.5\%) \\
\textbf{\ours} & \textbf{14} & \textbf{27 (3.8\%)} & 383 (53.2\%) & 97 (40.4\%) & 480 (50.0\%) \\
\bottomrule
\end{tabular*}

\medskip
\begin{tabular*}{\textwidth}{@{\extracolsep{\fill}}llrr@{}}
\toprule
Metric & Test unit & Difference [95\% CI] & $p$ \\
\midrule
Attack success & 42 $\to$ 14 families, 0/28 discordant & $-28.6$ [$-35.3, -21.9$] & $3.7\cdot10^{-9}$ \\
Benign utility & permutation over 48 families & $-4.5$ [$-9.4, -0.5$] & 0.044 \\
\bottomrule
\end{tabular*}
\end{table}


\paragraph{Block-all-sinks reference.} The ``Block all sinks'' rows in Tables~\ref{tab:test1} and~\ref{tab:test2}
estimate a policy that blocks every sink call, using \cNone{} trajectories. A benign case counts as successful if
\cNone{} succeeded without calling any sink. An attack case counts as harmful if \cNone{} caused harm through the
answer. On \tOne, at the same attack success rate (1.2\%), this policy keeps only 7.4\% of benign utility, compared with
38.6\% for \ours. The estimate is biased in two directions. For security, it favours the blocking policy, because a
blocked agent might move the confidential value into its answer. For utility, it disfavours the blocking policy,
because the agent could still finish the read-only part. The true utility gap may therefore be smaller.

\subsection{AgentDojo in English and Vietnamese}

% Nguồn duy nhất: facts.md, mục "AgentDojo: bảng đầy đủ mọi lượt". Sandwich có tấn công EN+VI chạy 06--07/11 (WP6),
% không tấn công 14/11 (C11); BrokerGuard 13--14/11 (C10).
% In đậm: tiện ích cao nhất và số tấn công thấp nhất trong mỗi bộ và cột (bỏ qua khi các ô bằng nhau).
\begin{table}[!htbp]
\centering
\caption{AgentDojo v1.2.2, \code{important\_instructions} attack. Each cell: utility under attack / successful
attacks (utility without attack). ``--'': not run. $^\dagger$Includes 2 cases in which every attack call was
blocked.}\label{tab:agentdojo}
\footnotesize
\setlength{\tabcolsep}{3pt}
\begin{tabular*}{\textwidth}{@{\extracolsep{\fill}}llrrrr@{}}
\toprule
 & & \multicolumn{2}{c}{Llama-3.1-8B} & \multicolumn{2}{c}{Gemini 3 Flash} \\
\cmidrule(lr){3-4}\cmidrule(l){5-6}
Suite (pairs, tasks) & Configuration & EN & VI & EN & VI \\
\midrule
Banking (144, 16) & \cNone & 58 / 17 (7) & 44 / 6 (6) & 93 / 4 (13) & 99 / 0 (12) \\
 & \cSand & 62 / 10 (7) & 45 / 6 (5) & -- & -- \\
 & \ours & \textbf{72} / \textbf{3} (8) & \textbf{54} / \textbf{1} (5) & -- & -- \\
\midrule
Slack (105, 21) & \cNone & 28 / 16 (7) & \textbf{26} / 0 (7) & 62 / 1 (18) & 59 / 0 (16) \\
 & \cSand & \textbf{35} / 10 (8) & 19 / 0 (6) & -- & -- \\
 & \ours & 15 / \textbf{6}$^\dagger$ (5) & 13 / 0 (5) & -- & -- \\
\bottomrule
\end{tabular*}
\end{table}


Table~\ref{tab:agentdojo} reports the AgentDojo results. Without defence, Llama is successfully attacked 17/144 times
on Banking and 16/105 times on Slack in English, but only 6 and 0 times in Vietnamese. The lower count in Vietnamese
does not mean that Vietnamese is safer. Llama's utility also drops (from 58 to 44 on Banking), and in Vietnamese Slack
two of the five injection tasks cannot be solved by the reference solution. \cSand{} reduces successful attacks to 10
on both English suites but makes no difference in Vietnamese. Without attacks, its utility differs from \cNone{} by at
most one task.

\ours{} reduces successful attacks further, to 3 on Banking and 6 on Slack in English, and 1 and 0 in Vietnamese. Two
of the six Slack cases are counted even though every attack call was blocked, because AgentDojo's scoring function
relies on the calls the model generated. Counting only executed calls, the figure is 4/105. The remaining eight cases
across the two suites and two languages are password changes (4 cases, through a sink without a destination), opening a
malicious URL with a read tool (3 cases) and sending a link to a person the user had named (1 case). On Slack the
destination is a recipient name, so name matching has little discriminating power. With real tools, the broker should
resolve names through a host-side directory. Utility rises on Banking (72 versus 62 for \cSand) but falls sharply on
Slack (15 versus 35), because many Slack tasks take their destination from a web page or a message that the user did not
write. Even without attacks, 9 of the 21 English Slack tasks have calls blocked by \rDest. \rSig{} never fires in
English and fires 14 times on Vietnamese Slack.

\subsection{Adaptive Attacks}\label{sec:adaptive}

The authors wrote 48 adaptive attack cases on Dev families, based on feedback from defended runs on Dev. Under
\cNone, this set is as strong as the corresponding static cases: 26, 15 and 24 harmful families for Gemma, Llama and
Qwen3.5, compared with 29, 12 and 22 for standard Vietnamese. Under \ours, these drop to 1, 0 and 0. This is only one
round of manual adaptation, not a white-box attack generated from the rules. The set does not target the gaps
described in Section~\ref{sec:discussion}.
